\documentclass[11pt, aspectratio=1610, xcolor={dvipsnames}]{beamer}

% font packages
\usepackage[utf8]{inputenc}
\usepackage[T1]{fontenc}
\usepackage{lmodern}

% misc packages
\usepackage{booktabs}
\usepackage{mathtools}
\usepackage{multirow}
\usepackage[round]{natbib}
\usepackage{transparent}

% theme and colours
\usetheme{Montpellier}
\setbeamercolor*{structure}{bg=PineGreen!10,fg=PineGreen}
\setbeamercolor*{footline}{bg=PineGreen!10,fg=PineGreen}

% line spread
\linespread{1.5}

% paths
\graphicspath{{../thesis/} {../../stata/figures/pdf_tex/}}

\makeatletter
\def\input@path{{../../stata/figures/pdf_tex/}}
\makeatother

\pdfsuppresswarningpagegroup=1

% set apa style
\bibliographystyle{apalike}


\title[Putting the GDR's Legacy Effect under the Microscope]{\texorpdfstring{Putting the GDR's Legacy Effect \linebreak under the Microscope:}{Putting the GDR's Legacy Effect under the Microscope:}}
\subtitle{Eastern Female STEM Professionals in Reunified Germany.}

%\author{\href{mailto:johannes.schuett@fu-berlin.de}{Johannes Schütt}}
\author{\texorpdfstring{\href{mailto:johannes.schuett@fu-berlin.de}{Johannes Schütt}}{Johannes Schütt}}

\institute{\texorpdfstring{Free University of Berlin \linebreak M.Sc. Public Economics \linebreak\linebreak Supervisor: Prof. Natalia Danzer, Ph.D.}{}}

\date{4 July 2024}

% logo on cover page
\titlegraphic{
	\vspace{5mm}
	\includegraphics[width=3cm]{fu_logo.pdf}
}

% define custom footline
\newcommand{\Ffootline}{%
	\hfill
	\insertframenumber/\inserttotalframenumber % The center
	\hfill
}

% set footline
\setbeamertemplate{navigation symbols}{}
\setbeamertemplate{footline}{%
	\ifnum\thepage=1
	\begin{beamercolorbox}[wd=\paperwidth,ht=0ex,dp=1ex]{footlinecolor}
	\end{beamercolorbox}
	\else
	\begin{beamercolorbox}[wd=\paperwidth,ht=0ex,dp=1ex]{footlinecolor}%
		\tiny\hspace*{4mm} \Ffootline \hspace*{4mm}
	\end{beamercolorbox}
	\fi
}

% bottom right transparent logo
\logo{
	{\transparent{0.6}
		\includegraphics[width = 1cm]{fu_seal.pdf}
		\hspace{1mm}
		\vspace{1mm}
	}
}

% sectiontitleframes
\AtBeginSection[]{
	\begin{frame}
		\vfill
		\centering
		\begin{beamercolorbox}[sep=8pt,center,shadow=false,rounded=false]{title}
			\usebeamerfont{title}\insertsectionhead\par%
		\end{beamercolorbox}
		\vfill
	\end{frame}
}

% special fonts
\newcommand{\highlight}[1]{\textbf{\textcolor{PineGreen}{#1}}}
\newcommand{\todo}[1]{\textbf{\textcolor{red}{#1}}}

% frame continuation for references
\makeatletter
\setbeamertemplate{frametitle continuation}{\gdef\beamer@frametitle{}}
\makeatother

% table def
\def\sym#1{\ifmmode^{#1}\else\(^{#1}\)\fi}

% change alignment for single cell inside tabular
\newcommand*{\myalign}[2]{\multicolumn{1}{#1}{#2}}


\begin{document}
	
	\begin{frame}[plain]
		\maketitle
	\end{frame}
	
	\begin{frame}
		\frametitle{}
		{\linespread{1}
		\tableofcontents
		}
	\end{frame}
	
	\section{Motivation}
	\begin{frame}
		\frametitle{}
		
		\begin{itemize}
			\item How do social norms influence our occupational decisions?
			\item Do social norms persist when the institutions that uphold them vanish?
			\item Consider these questions with regard to Eastern females who work in traditionally male-dominated work sectors after the socialist GDR collapses.
		\end{itemize}
		
	\end{frame}
	
	\subsection{Previous Literature}
	\begin{frame}
		\frametitle{Previous Literature}
		
		\begin{itemize}
			\item Econometric literature describes \highlight{persistence of East German social norms and culture} after Reunification.
			\item \cite{Lippmann2018} use PISA data from \highlight{2003} to show that the \highlight{performance gap of girls in maths} in the former GDR is being reduced.
			\item \cite{Jessen2023} find that an above-median \highlight{inflow of East German migrants} into West Germany led to \highlight{increase in labour supply} among employed \highlight{West German females}.
		\end{itemize}
		
	\end{frame}
	
	\begin{frame}
		\frametitle{}
		
		\begin{itemize}
			\item \cite{FriedmanSokuler2020} investigate \highlight{vertical persistence} (across generations) and \highlight{horizontal diffusion} (through peers) of socialist norms.
			\item Jewish immigrants from the former \highlight{Soviet Union to Israel} in the \highlight{early 1990s}.
			\item Females from FSU tend to \highlight{avoid work in the educational or social sector} and instead opt for \highlight{STEM} fields.
		\end{itemize}
		
	\end{frame}
	
	\begin{frame}
		\frametitle{}
		
		\textcolor{PineGreen}{\underline{Definition.}} STEM, field and curriculum centred on education in the disciplines of \textbf{S}cience, \textbf{T}echnology, \textbf{E}ngineering, and \textbf{M}athematics \textcolor{darkgray}{\citep{Hallinen2024}}.
		
	\end{frame}
	
	\begin{frame}
		\frametitle{}
		
		\begin{itemize}
			\item Even females who \highlight{immigrated as infants} are more likely to pursue STEM subjects in education than native Israeli females \highlight{(vertical persistence)}.
			\item \highlight{Native Israeli females} who attend school with high concentration of FSU immigrants shift their educational choices towards STEM \highlight{(horizontal diffusion)}.
		\end{itemize}
		
	\end{frame}
	
	\subsection{Developing a Research Question}
	\begin{frame}
		\frametitle{Developing a Research Question}
		
		\textcolor{PineGreen}{\underline{Initial Idea.}} Do Eastern Female STEM professionals have an impact on Western Females after the Reunification?
		
		\vspace{1cm}
		
		General trend in Germany regarding STEM professions using the SOEP…
		
	\end{frame}
	
	\begin{frame}
		\frametitle{}
		
		\begin{itemize}
			\item \highlight{Share of STEM professionals} within demographic groups in Germany over time.
			\item Group by: \highlight{East origin} (location in 1989) and \highlight{gender}.
			\item Observe the years 1990 (year of Reunification) -- 2017.
			\item All individuals in \highlight{working age} (17--65).
			\item Both employed and unemployed.
			\item Apply the SOEP's individual weights.
		\end{itemize}
		
	\end{frame}
	
	\begin{frame}
		\frametitle{}
		
		{\linespread{1}\small
			\begin{table}[h]
				\centering
				\caption{Examples of STEM Professions (sorted by Frequency)}
				\label{tab:stem_examples}
				
				\begin{tabular}{ll}
					\toprule
					ISCO-88 Code & Professional title\\
					\midrule
					{[2149]} & Architects, engineers and related scientists\\
					{[3152]} & Health, environmental and quality inspectors\\
					{[3119]} & Materials and engineering specialists\\
					{[2142]} & Civil engineers\\
					{[3121]} & Data processing assistants\\
					{[3111]} & Chemo- and physicotechnician\\
					{[2145]} & Mechanical engineers\\
					{[3114]} & Electronics and telecommunications technology\\
					{[3120]} & Data processing specialists\\
					{[2144]} & Electronics and telecommunications engineers\\
					… &\\
					\bottomrule
				\end{tabular}
			\end{table}
		}
		
	\end{frame}
	
	\begin{frame}
		\frametitle{}
		
		\begin{figure}[h]
			\centering
			\caption{Time Trend (by Region and Gender), 1984--2017}
			\label{fig:trend}
			\resizebox{75mm}{!}{\input{trend.pdf_tex}}
		\end{figure}
		
		{\scriptsize
			\textcolor{darkgray}{Source: \cite{SOEP2023}. Own calculations.}
		}
		
	\end{frame}
	
	\begin{frame}
		\frametitle{}
		
		\begin{figure}[h]
			\centering
			\caption{Time Trend (by Region and Gender), 1984--2017}
			\label{fig:trend_highlight}
			\resizebox{75mm}{!}{\input{trend_highlight.pdf_tex}}
		\end{figure}
		
		{\scriptsize
			\textcolor{darkgray}{Source: \cite{SOEP2023}. Own calculations.}
		}
		
	\end{frame}
	
	\begin{frame}
		\frametitle{}
		
		\begin{figure}[h]
			\centering
			\caption{Time Trend (by Region, Females only), 1990--1999}
			\label{fig:trend_zoomed}
			\resizebox{75mm}{!}{\input{trend_zoomed.pdf_tex}}
		\end{figure}
		
		{\scriptsize
			\textcolor{darkgray}{Source: \cite{SOEP2023}. Own calculations.}
		}
		
	\end{frame}
	
	\begin{frame}
		
		\textcolor{PineGreen}{\underline{Research Question.}} What drives the downward trend of East German females in STEM professions in the first decade after Reunification?
		
	\end{frame}
	
	\subsection{Historical Background}
	\begin{frame}
		\frametitle{Historical Background}
		
		\begin{itemize}
			\item 92\% of all \highlight{female apprentices} in the GDR were concentrated in thirteen \highlight{occupational groups}, compared to five in the FRG \textcolor{darkgray}{\citep{Menschik1974}}.
			\item \highlight{Horizontal segregation} in the labour sector in the GDR was not nearly as high as in the societies of the Western bloc \textcolor{darkgray}{\citep{Lane1983}}.
		\end{itemize}
		
	\end{frame}
	
	\begin{frame}
		\frametitle{}
		
		\begin{itemize}
			\item Of the \highlight{female apprentices} who successfully passed their skilled labour examination (\emph{Facharbeiterprüfung}) in 1980, \highlight{17\%} were in the \highlight{STEM field} \textcolor{darkgray}{\citep[p. 293]{DDRJahrbuch1981}}.
			\item Unlike in the FRG, in the GDR there was the \highlight{same school curriculum for boys and girls} with the same \highlight{emphasis on science and maths} \textcolor{darkgray}{\citep{FuchsSchuendeln2016, Lippmann2018}}. 
		\end{itemize}
		
	\end{frame}
	
	\begin{frame}
		\frametitle{}
		
		\begin{itemize}
			\item In \highlight{1949}, females made up \highlight{40\% of the total labour force} both in the GDR and in the FRG.
			\item Unlike the West German females, \highlight{East German females remained in the labour market} after the economic upturn in the 1960s \textcolor{darkgray}{\citep{Menschik1974}}.
		\end{itemize}
		
		%\todo{Notizen einarbeiten!}
		
	\end{frame}
	
	\section{Empirical Approach}
	\begin{frame}
		\frametitle{}
		
		\begin{itemize}
			\item Compare \highlight{East German females} with  \highlight{East German males}.
			\begin{itemize}
				\item East German males and females had the same probability\linebreak of working as a STEM professional in 1990 (what about before?).
				\item Examine in what way they each were affected by the Reunification.
			\end{itemize}
			\item Regressor of interest: \highlight{years after Reunification $\times$ being female}.
			\item Use suitable controls to \highlight{explain what drives the downward slope}.
		\end{itemize}
		
		
	\end{frame}
	
	\subsection{Econometric Specification}
	\begin{frame}
		\frametitle{Econometric Specification}
		
			{\linespread{1}\normalsize
				\begin{align}
					&\begin{aligned}
						\label{eq:spec}
						STEM_{itr} =& \beta_{0} + \beta_{1} Dist\_Reunification_{t} + \beta_{2} Female_{i} +\\
						&\beta_{3} Dist\_Reunification \times Female_{it} + X_{itr} \gamma + \epsilon_{itr}
					\end{aligned}
				\end{align}
				
				\vspace{1cm}
				
				\hspace{6.2em} with $Dist\_Reunification_{t} = Survey\_Year_{t} - 1990$.
			}
		
	\end{frame}
	
	\subsection{Data}
	\begin{frame}
		\frametitle{Data}
	\end{frame}
	
	\section{Preliminary Results}
	\begin{frame}
		\frametitle{}
	\end{frame}
	
	\subsection{Main Findings}
	\begin{frame}
		\frametitle{Main Findings}
		
		{\linespread{1}\tiny
			\begin{table}[h]
				\centering
				\caption{Margins (dy/dx) of Cloglog Regressions with y = being a STEM Professional}
				\label{tab:margins}
				\begin{tabular}{l*{4}{c}}
					\toprule
					&\multicolumn{1}{c}{(1)} & \multicolumn{1}{c}{(2)}         &\multicolumn{1}{c}{(3)} & \multicolumn{1}{c}{(4)} \\
					\midrule
					Years after Reunification & -0.0008 & -0.0027 & -0.0027 & -0.0027 \\
					& (0.0016) & (0.0017) & (0.0017) & (0.0017) \\
					\addlinespace
					Female & -0.0022  & -0.0032 & -0.0032 & -0.0032 \\
					& (0.0088) & (0.0088) & (0.0088) & (0.0088) \\
					\addlinespace
					Years after Reunification $\times$ Female &     -0.0121\sym{***}&     -0.0101\sym{***}&     -0.0101\sym{***} &     -0.0101\sym{***}\\
					&    (0.0028)         &    (0.0029)         &    (0.0029)         &    (0.0029) \\
					\addlinespace
					Residence in West Germany&                     &      0.0474\sym{***}&      0.0474\sym{***}&      0.0474\sym{***}\\
					&                     &    (0.0174)         &    (0.0174)         &    (0.0174)         \\
					\addlinespace
					Residence in West Germany $\times$ Female         &                     &     -0.0517\sym{*} & -0.0517\sym{*} & -0.0517\sym{*} \\
					& &    (0.0309)         &    (0.0309)         &    (0.0309) \\
					\midrule
					Observations        &       14,632         &       14,632         &       14,632         &       14,632         \\
					\bottomrule
				\end{tabular}
			\end{table}
		}
		
		{\scriptsize
			\textcolor{darkgray}{Source: \cite{SOEP2023}. Own calculations.}
		}
		
	\end{frame}
	
	\begin{frame}
		\frametitle{}
		
		\begin{itemize}
			\item Effect driven by Attrition?
			\item New panel entrants?
		\end{itemize}
		
	\end{frame}

	\begin{frame}
		\frametitle{}
		
		\begin{figure}[h]
			\centering
			\caption{Net Switches among Eastern Females, 1990--1999}
			\label{fig:eastern_female_tracking}
			\resizebox{75mm}{!}{\input{eastern_female_tracking.pdf_tex}}
		\end{figure}
		
		{\scriptsize
			\textcolor{darkgray}{Source: \cite{SOEP2023}. Own calculations.}
		}
		
	\end{frame}
		
	\subsection{Discussion}
	\begin{frame}
		\frametitle{Discussion}
	\end{frame}
	
	\section{Conclusion}
	\begin{frame}
		\frametitle{}
	\end{frame}
	
	\section*{References}
	\begin{frame}[allowframebreaks]
		\frametitle{References}
		
		{\scriptsize
		\bibliography{../references.bib}
		}
		
	\end{frame}
	
	\appendix
	
	\section{Appendix}
	
	\subsection{Validity}
	\begin{frame}[plain]
		\frametitle{Validity}
	\end{frame}
	
	\begin{frame}
		\frametitle{}
		
		\begin{figure}[h]
			\centering
			\caption{Cohorts (Start of Occupation by Gender), 1945--1990}
			\label{fig:validity}
			\resizebox{75mm}{!}{\input{validity.pdf_tex}}
		\end{figure}
		
		{\scriptsize
			\textcolor{darkgray}{Source: \cite{Mayer1995}. Own calculations.}
		}
		
	\end{frame}
	
	\begin{frame}
		\frametitle{}
		
		{\linespread{1}\tiny
			\begin{table}[h]
				\centering
				\caption{Share of Females in the Skilled Labour Force by Economic Sector in 1971 (in \%)}
				\label{tab:gdr_yearbook}
				
				\begin{tabular}{llrrrrrrr}
					\toprule
					Labour sector & Selected Segment & \multicolumn{5}{c}{Below the age of} & & \multirow[c]{2}{*}{Total} \\
					& & \myalign{c}{25} & \myalign{c}{30} & \myalign{c}{40} & \myalign{c}{50} & \myalign{c}{60} & \myalign{c}{60\texttt{+}} & \\
					\midrule
					Technical sciences & & 32.6 & 16.4 & 8.6 & 4.3 & 3.0 & 1.1 & 9.3 \\
					& Mechanical engineering & 24.2 & 10.1 & 4.1 & 1.1 & 0.5 & 0.2 & 4.8 \\
					& Textile technology (mechanical) & 84.0 & 67.8 & 41.3 & 23.3 & 16.8 & 5.2 & 36.7 \\
					& Chemical engineering & 60.3 & 50.3 & 39.1 & 38.0 & 29.2 & 17.2 & 42.1 \\
					& Automation engineering & 18.3 & 6.1 & 3.8 & 0.6 & 1.4 & 0.0 & 5.8 \\
					& Electrical engineering & 13.6 & 4.1 & 2.7 & 1.2 & 0.3 & 0.1 & 2.7 \\
					& Energy technology & 16.3 & 8.9 & 4.6 & 2.9 & 2.2 & 0.9 & 6.1 \\
					& Construction industry technology & 34.3 & 17.6 & 6.3 & 1.6 & 0.5 & 0.2 & 6.6 \\
					& Mining engineering &  15.8 & 3.5 & 1.4 & 0.5 & 0.3 & 0.0 & 1.5 \\
					Economic sciences & & 72.4 & 53.2 & 37.5 & 18.8 & 15.0 & 8.9 & 30.4 \\
					Medicine, Agricultural sciences & & 53.4 & 40.3 & 29.8 & 14.7 & 10.2 & 4.1 & 25.3 \\
					& Medical and Pharmacy technology & 96.7 & 89.9 & 84.2 &  64.8 & 36.6 & 20.1 & 76.8 \\
					& Agricultural sciences & 46.0 & 29.3 & 21.6 & 10.3 & 6.7 & 1.3 & 18.3 \\
					Cultural, Educational, and Sports sciences & & 95.4 & 86.8 & 80.0 & 72.1 & 64.9 & 43.3 & 80.0 \\
					Literature, Journalism & & 71.4 & 46.8 & 35.6 & 34.6 & 32.7 & 17.3 & 34.9 \\
					\bottomrule
				\end{tabular}
			\end{table}
		}
		
		{\scriptsize
			\textcolor{darkgray}{Source: \cite[p. 442]{DDRJahrbuch1973}. Own calculations.}
		}
		
		
	\end{frame}
	
\end{document}
